\chapter{检索服务与联想结构}

\section{服务视图}

权威认知对象可以派生多个服务视图，包括精确索引、词项文档、稠密向量、实体与标签倒排、语义片段、召回线索、联想拓扑和其他读优化表示。服务视图可以独立失效、删除和重建，不改变权威对象身份。

稠密向量视图绑定完整的向量空间身份。向量空间身份至少能够区分模型或权重、任务、分词和预处理、池化、归一化方式及维度。不同向量空间不得直接混用向量分数。

\section{服务代次}

每类服务投影按照不可变代次发布。代次记录输入权威水位、实现身份、关键配置摘要、制品身份和必要的校验信息。构建完成并通过一致性检查后再切换当前代次。

查询返回前仍以权威对象为最终依据。服务代次负责高效召回和排序，不拥有认知真值。

\section{联想拓扑}

联想拓扑（AssociativeTopology）是从记忆、标签、经历片段、认知图式、社会关系和关联证据派生的可重建服务图。节点和边的服务权重由版本化构建规则生成。

拓扑构建遵守有限资源原则：出边质量具有固定或归一化预算；高连接度节点受到枢纽或扇出抑制；查询扩展具有跳数、节点、状态、时间和计算预算；共同来源祖先在需要时去重；被截断或未探索的质量能够进入诊断。

负向、矛盾和反例关系具有独立语义，不能在没有定义转换规则时直接作为普通非负传播边使用。

\section{关联证据与服务权重}

关联证据保存有来源、有类型的认知关系依据。服务层根据当前有效关联证据、关系类型、方向、极性、来源独立性、对象生命周期和查询意图计算拓扑边权。

系统不保存一个由检索次数不断累加的通用“关联强度真值”。被召回、排序或呈现不会自动生成新的关联证据。只有具有独立语义的外部观察、明确使用结果或领域修订才能成为新的长期支持来源。
